\chapter{The Store Is the Denotation}
\label{ch:rspace}

\section{Two senses of ``runs''}

Everything in the Machinery part is a construction on a language presentation. The
reader has been told, repeatedly, that the constructions are not metaphors: that a
cost-accounted theory meters, that a graded clause resolves a race, that a $\dok$
intervenes. Each of those verbs has two senses. In the first, ``the theory meters'' means
that the rewrite relation, read as a mathematical object, has a component that decreases.
In the second, it means that a machine somewhere is doing the subtracting, and that if one
pulled the plug the subtracting would stop. The Pledge promised the second sense when it
said that \emph{the network is the computer} and pointed at the Groq and Cerebras dies
and the FPGAs whose fabric comes and goes. This part is where the second sense is
delivered, and it is delivered in two installments of very different maturity.

The first installment is RSpace. It is a data structure, an algorithm over it, and a
production store with several years of operation behind it, and this chapter's claim is
that the store \emph{is the denotation}: that the compilation of rho into RSpace states is,
term for term, a denotational semantics, with the structural laws of the calculus holding
as literal equalities of hashtables. The second installment, in the next chapter, is a
compiler from rho through a name-free combinator calculus to a matrix machine that runs on
a graphics processor. It exists, it passes its tests, and it has not yet been measured on
anything that matters; we will say exactly how nascent it is. The two are realizations of
the same COMM rule on two hardware philosophies that reward opposite things, and the third
chapter of the part asks what that opposition prices.

A word about the voice. The Machinery part spoke of theories; this part speaks of a
particular implementation, f1r3node, written in Rust, and of a particular compiler,
Saph1r3. Where the book says what the code does, it has been checked against the code and
its documentation at the time of writing; where it says what the code's authors intend,
it says so. The book does not become a manual here. It becomes a report on whether the
abstractions of the preceding part are abstractions \emph{of} anything, and the answer is
that one of them has been running in production for some time and the other was compiled
for the first time this autumn.

\section{The store}
\label{rs:sec:store}

Modulo structural congruence the parallel fragment of rho, $(\Par, \nil)$, is the free
commutative monoid, and the free commutative monoid on a set is the finite-multiset
construction on it. Interpreting $\Par$ as multiset union therefore makes the denotation a
monoid homomorphism \emph{on the nose}, and structural congruence of the parallel skeleton
becomes literal equality of multisets --- the tightest fit syntax and semantics can have.
Bucketing the multiset by channel turns it into a hashtable. That hashtable is RSpace.

\begin{definition}[Tables]
\label{rs:def:table}
A \emph{table} is a finite partial map from names to \emph{polarity-homogeneous}
multisets of payloads: a key holds either a bag of output payloads
$x \mapsto \{\!|\,Q_1,\dots,Q_m\,|\!\}^{+}$ --- \emph{produces}, data --- or a bag of
input continuations $x \mapsto \{\!|\,(y_1)P_1,\dots,(y_n)P_n\,|\!\}^{-}$ --- \emph{consumes},
abstractions --- never a mixture. The keys are names, hence quoted processes: RSpace is
\emph{content-addressed}, and a channel's identity is the hash of the process it quotes.
\end{definition}

\begin{definition}[The RSpace denotation]
\label{rs:def:rspace}
\begin{align*}
  \sem{\nil}&=\varnothing, &
  \sem{x!(Q)}&=\{\,\sem{x}\mapsto\{\!|\sem Q|\!\}^{+}\,\},\\
  \sem{P\Par Q}&=\sem P \uplus_{k} \sem Q, &
  \sem{\fork(y\leftarrow x)P}&=\{\,\sem{x}\mapsto\{\!|(y)\sem P|\!\}^{-}\,\},\\
  \sem{*x} &= *\sem{x}, & \sem{@P} &= @\sem{P},
\end{align*}
where $\uplus_k$ is \emph{key-wise} multiset union: the bags $P$ and $Q$ install at a key
are united. Each prefix installs a singleton bag; $\Par$ co-locates produces and consumes.
A key whose two operands carry the same polarity stays homogeneous; a key receiving both
polarities is a \emph{redex}.
\end{definition}

\begin{proposition}[Exact homomorphism]
\label{rs:prop:hom}
$\sem{-}$ is a homomorphism of commutative monoids: $\sem\nil=\varnothing$ and
$\sem{P\Par Q}=\sem P\uplus_k\sem Q$. Hence the parallel laws of structural congruence
--- associativity, commutativity, unit --- hold as literal equalities of tables, and
$P\equiv Q \Rightarrow \sem P=\sem Q$.
\end{proposition}

\begin{proof}
$\uplus_k$ inherits the commutative-monoid laws of multiset union key-wise, and the empty
table is the unit. The prefix clauses are singletons, so reassociating or commuting a
parallel composition reassociates or commutes the corresponding bag unions, which multiset
equality absorbs. Multiset union has no idempotence to defeat: $a!(0) \Par a!(0)$ installs a
two-element bag at $a$, distinct from the one-element bag of $a!(0)$. Counts are kept.
\end{proof}

The reflection law $*@P \equiv P$ is validated separately, and the source for this
section~\cite{qcs} derives it rather than assuming it: in a universe stratified into two
sorts of set and two sorts of atom, with an involution swapping the colors, quotation is a
seal followed by a swap, dereference undoes it, and $*@P \equiv P$ is the statement that
the swap is its own inverse. The book does not need the stratification and does not rehearse
it. What it needs is that every generator of structural congruence holds in the table
model --- the parallel laws on the nose and reflection by the involution --- so that a
table is a well-defined function of a process modulo $\equiv$.

\begin{definition}[Eager reaction: the zip]
\label{rs:def:zip}
A key $x$ carrying both a produce-bag $\{\!|Q_1,\dots,Q_m|\!\}^{+}$ and a consume-bag
$\{\!|(y_1)P_1,\dots,(y_n)P_n|\!\}^{-}$ is a redex. Eager reaction fires it: choose
underlying lists of the two bags and \emph{zip} them; each matched pair reacts by
application, $(y_i)P_i \cdot Q_i = P_i\{@Q_i/y_i\}$, and the resulting process is merged
by $\uplus_k$ back into the table; the $|m-n|$ unmatched payloads of the longer bag remain
at $x$ as a homogeneous surplus. Its normal forms are the tables with no mixed-polarity key.
\end{definition}

Which $Q_i$ meets which $(y_j)P_j$, and which payloads survive as surplus, depend on the
chosen list orderings. This is rho's choice of which send meets which receive, now
localized at the key --- the \emph{resolver} of Chapter~\ref{ch:grading}, made into a
datum one can point at. The denotation $\sem{-}$ remains a function (plain $\uplus_k$, no
firing); reaction is the separate dynamics, a relation on tables confluent only up to the
equivalence that forgets the pairing. And the minimal context $[\,]\Par x!(Q)$ that
Chapter~\ref{sec:lts-hml} uses to label a transition is exactly the minimal store insertion
--- one produce at key $x$ --- that turns a pending consume into a redex. \emph{Context
labels are RSpace puts.}

\begin{remark}[Store monad plus zip]
\label{rs:rmk:bialgebra}
The table functor carries a monad structure --- the finite-multiset monad, lifted key-wise
--- and RSpace is this monad \emph{together with} eager reaction, presented as a
distributive law of the reaction over the store: a bialgebra in the sense of Turi and
Plotkin~\cite{turiplotkin}. This is the precise content of the description one hears of
RSpace as ``strictly more than a monad'': the monad is the store, the surplus is the zip.
Whether the bialgebra laws hold --- whether eager reaction is confluent up to resource
bisimilarity, whether the distributive law is one --- is filed in the source as an
obligation and not a theorem~\cite[\S8]{qcs}. The book reports it so.
\end{remark}

\section{Two layers, and what each is adequate for}
\label{rs:sec:layers}

Forget multiplicities and flatten a table into its set of context-labeled transitions ---
a consume at $x$ is an offer-input, a produce an offer-output, a mixed key an internal
redex --- and one has a natural transformation $\pi$ from tables to the behavior functor
of Chapter~\ref{sec:lts-hml}. The RSpace model sits \emph{above} the behavioral model and
projects onto it:
\[
  \text{Syntax} \xrightarrow{\ \sem{-}\ } \text{Tables (RSpace)}
                \xrightarrow{\ \pi\ } \text{Behaviors},
\]
and the composite is the context-behavior denotation the rest of the book works in.

\begin{definition}[Resource bisimilarity]
\label{rs:def:resbis}
Two tables are \emph{resource bisimilar} when they are bisimilar as coalgebras of the
table functor: same keys, same multiplicities at each key, and matching successors under
eager reaction. It refines context bisimilarity, with $\pi$ the quotient that forgets
multiplicity and keys.
\end{definition}

The source states two full-abstraction theorems~\cite{qcs}: at the resource layer,
resource bisimilarity coincides with equality in the final coalgebra of the table functor;
at the behavioral layer, context bisimilarity coincides with equality in the final
coalgebra of the behavior functor, and because the labels are contexts the behavioral
equivalence is a congruence by construction, so full abstraction holds for a congruence
with no closure step. Both proofs are sketches, and the obligations section of the source
lists what would discharge them: existence of the final coalgebras at the right size,
eager confluence, the bialgebra laws, and weak pullback preservation for the functors
involved. We state the theorems as the source states them and at the confidence the source
has.

Why two layers, and not one? Because they answer different questions. The resource layer
is \emph{the machine}: it knows that there are two messages at $a$ and not one, it knows
which key is which, and it is what a validator must reproduce to the byte when it replays a
block. The behavioral layer is \emph{what a program means}: it does not care whether the
two messages are one or two, only what the system can do in the presence of an observer.
That the first projects onto the second is the whole of the claim that running the program
on the store computes its meaning, and that the projection is a quotient --- forgetting
counts and keys --- is the precise sense in which the machine knows more than the meaning
does. Chapter~\ref{ch:grading}'s multiplicity factor, which counted candidates by occurrence
and not by term, was this quotient refused: a rate is a resource-layer quantity.

\begin{remark}[What the formal model does not cover]
\label{rs:rmk:gap}
The denotation above is for monadic single-channel input with no patterns. The store the
next section describes implements joins (a consume waiting on several channels at once),
spatial pattern matching on the received process, persistent produces and consumes, and
non-consuming reads. None of these is in the formal model; the source names joins as
future work. A reader should hold the two apart: the theorem is about a fragment, and the
implementation is of the whole language. That the fragment's theorem is the right shape
for the whole is likely and is not proved.
\end{remark}

\section{The store as built}
\label{rs:sec:built}

f1r3node's store is a Rust crate, \texttt{rspace++}, behind an interface generic in four
types: channels $C$, patterns $P$, data $A$ and continuations $K$. For rho, $C$ is a
process (a channel is a quoted term), $P$ is a binding pattern, $A$ is a list of processes
with a random-seed carrier, and $K$ is a tagged continuation. The interface has a
\emph{produce} --- put data on a channel, persistently or not --- and a \emph{consume} ---
wait on a \emph{vector} of channels with a vector of patterns and a continuation,
persistently or not, with a set of indices marked as non-consuming reads. Every construct of
Definition~\ref{rs:def:rspace} is one call: a send is a produce, a receive is a consume,
and the interpreter's reduction loop, on evaluating a send, calls produce, and on
evaluating a receive, calls consume, with the optional \texttt{where} clause passed as a
guard the matcher evaluates per candidate --- which is to say that the clause of
Chapter~\ref{ch:grading} already reaches the store.

\paragraph{Matching is the zip, with patterns.} When a produce arrives, the store looks up
the joins registered on that channel, then the waiting continuations; candidates are
shuffled with a random generator seeded from the deploy, the data on each joined channel is
matched speculatively against its pattern, and failed attempts are rolled back. A match
fires a COMM, feeding the continuation back to the interpreter with the matched data
substituted. With no match, the item is stored. This is Definition~\ref{rs:def:zip} with
two additions the formal model lacks --- the join, which makes a redex a tuple of keys
rather than one, and the pattern, which makes ``matched'' a structural question about the
payload rather than a mere co-location --- and one the formal model leaves open: the
shuffle, which is the resolver, and which is seeded so that it is \emph{reproducible}.

\paragraph{Hot store and history.} The working set is an in-memory map from channels to
their data, continuations and joins. A \emph{checkpoint} converts the working set's
mutations into updates of a 256-ary radix trie whose nodes are hashed with Blake2b-256, so
the trie is a Merkle trie and its root is a 32-byte summary of the entire state; the trie
lives in an embedded key-value store on disk, in three databases --- history nodes,
checkpoint roots, and cold leaf data. A checkpoint is a root hash together with the event
log since the last one. That is the history monad of Chapter~\ref{ch:histmonad} with a
particular representation: $\Hist$ said a state carries a word of events; the trie says
that the word may be summarized by a hash, that two states with the same hash are the same
state, and that the summary may be recomputed by anyone holding the leaves. The
representation is not incidental. It is what makes the next paragraph possible.

\paragraph{Replay.} A second implementation of the same interface, \texttt{ReplayRSpace},
takes a checkpoint and the recorded event log and \emph{forces} the recorded COMMs: when
the interpreter, re-running the same deploys from the same root, reaches a race, the replay
store consults the log and resolves it as it was resolved the first time. Non-deterministic
system processes --- calls to an external model, to a network service --- are replayed from
their recorded outputs rather than re-executed. The replayer reaches the same root hash or
it does not, and the comparison is the whole of what a validator checks. We will say more
about this in the third chapter, because it is where the execution model of the shard and
the three rungs of Chapter~\ref{ch:do} meet: the executing node resolves; the replaying
node observes; and a counterfactual, in Chapter~\ref{ch:do}'s sense, is a replay from a
pre-state with one record changed.

\paragraph{Charging.} A wrapper, \texttt{ChargingRSpace}, meters every store operation. It
pre-charges for storage when a produce or consume is stored, refunds the pre-charges when a
COMM fires and charges one combined COMM cost instead, and re-charges for persistent items
that remain. A field on the COMM event records whether a produce or a consume triggered it,
so that play and replay --- in which task scheduling may cause a different operation to
trigger the same COMM --- charge identically. This is the cost monad of
Chapter~\ref{ch:costmonad} and the metered calculus of Chapter~\ref{ch:carho} running, and
it is worth noticing which of their features the implementation had to get right before
the book had a name for it: the \emph{order-independence} of charging, which the book
derives from the requirement that a validator's bill agree with the proposer's, and which
Proposition~\ref{wt:prop:gate} states as the gate deciding support while the weight
distributes over it.

\paragraph{Merging.} A separate module analyzes event logs for conflicts: whether one
event depends on another, whether two events race on a non-persistent operation, and why.
This is what lets a block have more than one parent --- two validators' state changes
merged when their logs do not conflict --- and it is the first appearance in the
implementation of the contention relation Chapter~\ref{ch:grading} defined: two events
conflict when they share a datum or a receipt, and independent events commute.

\section{What the machine knows that the calculus does not}
\label{rs:sec:knows}

Three facts about the store are not facts about rho, and the book's honesty about the
relationship between the two requires naming them.

\emph{The store is locked, not lock-free.} The architecture document describes the store
as lock-free; the code uses 4096 pre-allocated mutex stripes selected by channel hash,
acquired in sorted order in two phases, with first-in-first-out waiters. Branches of a
parallel composition are evaluated as concurrent tasks, and the stripes serialize their
access to each channel. There is nothing wrong with this; it is how a general machine
realizes a tuple space. But it means the realization has a scheduler, and the scheduler is
not in the calculus. Determinism across nodes comes from three decisions, each outside
rho: sends are evaluated before receives, the candidate shuffle is seeded from the deploy,
and replay forces the log.

\emph{Some joins are blocked.} A consume waiting twice on the same channel is refused by the
implementation, because the candidate-extraction algorithm produces unmatched COMMs for
overlapping patterns on one channel; and the cost of some same-channel joins involving
persistent operations is acknowledged to be non-deterministic. These are documented
limitations, with issue numbers. The calculus has no such restriction. The gap is small and
it is real.

\emph{The bottleneck is not where one would guess.} Measured, checkpoint durability is about
one percent of the cost and spawning a replay runtime is under four milliseconds; the
sustained throughput limit is the block replay pipeline, and the delay from block inclusion
to finality is among its causes. Nothing in the repository relates the store to graphics
processors or other accelerators; a design document lists them as a non-goal. The matrix
note of the next chapter says why, from the other side: an RSpace-style interpreter does
approximately the opposite of everything an accelerator rewards.

That last sentence is the hinge of the part. The store is mature, correct in production,
and shaped for a general machine. It is a hashtable of variably shaped terms reached by
pointer, matched by a recursive walk, with conflicts resolved by locks. A graphics
processor wants flat indexable state, uniform work across lanes, bulk-synchronous phases,
and standard sparse primitives. Both machines run the same COMM rule. The question of the
next chapter is what a compiler standing between them has to do.

\begin{remark}[What this chapter claims, and at what confidence]
\label{rs:rmk:claims}
\emph{Implemented and operating}: every send is a produce and every receive a consume on a
content-addressed store; channels are quoted processes; a rendezvous is literally the COMM
rule; the whole state is summarized by one Merkle root; a validator replays the log and
checks the root. \emph{Formally stated, with short proofs}: the denotation is an exact
monoid homomorphism, and reflection holds in the model. \emph{Formally stated, proofs
sketched, obligations filed}: full abstraction at both layers; the bialgebra laws; eager
confluence. \emph{Not modeled}: joins, patterns, persistence, peeks. The chapter has tried
to keep these four registers apart on every page.
\end{remark}

% ===================================================================
